\chapter*{Part VII Guide: the Checklist from an InfoEdge Data Scientist}
\addcontentsline{toc}{chapter}{Part VII Guide: Checklist Coverage Map}
\stuck{InfoEdge DS interview experience (full video)}{\ytid{nWyBQhGnHaM}}

A candidate who secured an InfoEdge data scientist offer shared the checklist below of what to know for the interviews. This page maps every item to
where it is taught. Items that the earlier chapters already covered fully point there; items that were missing or only touched on are taught in the
new chapters 31--40, which follow the same pattern as the rest of the book: concepts from zero, many small examples, a ``where this is used in ML''
section, and interview questions with full answers at the end. Probability, statistics and linear algebra are also covered at OA depth in your
separate Maths book; here they are taught again with an interview and ML-application focus.

\begin{center}\small
\begin{tabularx}{\linewidth}{>{\raggedright\arraybackslash}X >{\raggedright\arraybackslash}p{4.6cm}}
\toprule
\textbf{Checklist item} & \textbf{Where}\\
\midrule
\multicolumn{2}{l}{\textbf{\color{accent}EDA}}\\
Reason behind every step; data visualisations and how to read each graph & Ch.\ 31 (new), Ch.\ 26\\
Univariate statistics: mean, median, mode, range, standard deviation, variance & Ch.\ 31 (new)\\
Missing values: MCAR/MAR/MNAR, dropping, mean/median/mode imputation pros and cons & Ch.\ 13, Ch.\ 33 (new)\\
ML models to predict missing values (KNN, MICE, MissForest, XGBoost imputer) & Ch.\ 33 (new), Ch.\ 13\\
Leaky variables & Ch.\ 1, Ch.\ 27, Ch.\ 33 (new)\\
Outliers: z-score, 68--95--99.7 rule, IQR formula, box plot interpretation & Ch.\ 33 (new), Ch.\ 13\\
Scaling: standardisation vs normalisation, min--max, robust & Ch.\ 34 (new)\\
Numerical to categorical: binning, clustering & Ch.\ 34 (new)\\
Categorical to numerical: ordinal, one-hot, target, frequency encoding; dummy variable trap & Ch.\ 34 (new)\\
Imbalanced data: under/oversampling, weighting, NearMiss, random deletion, SMOTE, Borderline-SMOTE, ADASYN & Ch.\ 13, Ch.\ 34 (new)\\
Bivariate: Pearson, Spearman (ranks), point-biserial, phi; correlation matrix & Ch.\ 32 (new)\\
Multivariate: multicollinearity, VIF in terms of $R$ & Ch.\ 32 (new), Ch.\ 2\\
Skewed data: log, power, Box--Cox transforms & Ch.\ 34 (new), Ch.\ 31\\
\midrule
\multicolumn{2}{l}{\textbf{\color{accent}Basic ML}}\\
OLS, 5 assumptions (detect and fix), MLE for linear regression, loss & Ch.\ 2, Ch.\ 12\\
L1, L2, elastic net; Huber; homo-/heteroscedasticity & Ch.\ 1, Ch.\ 2, Ch.\ 4\\
GLM maths; Poisson and logistic as GLMs & Ch.\ 4, Ch.\ 3\\
SVM, hinge, kernels (polynomial, Gaussian/RBF), Mercer, SVR, margin, hyperplane & Ch.\ 9, Ch.\ 4\\
Naive Bayes (3 types), Laplace smoothing, independence assumption & Ch.\ 10\\
K-Means and variants (k-medoids, k-modes, k-prototypes, k-means++), soft vs hard clustering & Ch.\ 15, Ch.\ 35 (new)\\
GMM, EM, DBSCAN, Swiss roll & Ch.\ 15, Ch.\ 35 (new)\\
PCA (scree, explained variance, eigen intuition), SVD, Fisher LDA & Ch.\ 14, Ch.\ 10, Ch.\ 40 (new)\\
t-SNE, SNE, perplexity, Student-t; linear vs non-linear and deterministic vs non-deterministic DR & Ch.\ 14, Ch.\ 35 (new)\\
Decision trees, Gini, entropy, information gain, Extra Trees, RF, bagging, boosting, AdaBoost vs GB & Ch.\ 6--8\\
Bias vs variance with graphs; tackling under- and overfitting & Ch.\ 1\\
\midrule
\multicolumn{2}{l}{\textbf{\color{accent}Deep Learning}}\\
ANN, perceptron structure, activations, training, backprop & Ch.\ 18, Ch.\ 36 (new: perceptron algorithm)\\
GD variants (stochastic, batch, mini-batch), momentum, Adam; loss functions & Ch.\ 5, Ch.\ 18\\
Transfer learning, CNN layers, $\frac{n+2p-f}{s}+1$ & Ch.\ 19\\
RNN, GRU, LSTM & Ch.\ 20\\
Basics of CV, NLP, reinforcement learning, federated learning & Ch.\ 36 (new), Ch.\ 16--17\\
Time series: ANOVA, ARIMA, SARIMAX & Ch.\ 37 (new), Ch.\ 20\\
\midrule
\multicolumn{2}{l}{\textbf{\color{accent}Statistics, Probability, Linear Algebra}}\\
Descriptive vs inferential, population vs sample, CLT, p-value, hypothesis testing & Ch.\ 39 (new)\\
Distributions with real-life examples; PDF, PMF; expectation, variance, moments & Ch.\ 38 (new)\\
Skewness and kurtosis and their types & Ch.\ 31 (new), Ch.\ 38 (new)\\
Eigenvectors, eigenvalues, rank & Ch.\ 40 (new)\\
\bottomrule
\end{tabularx}
\end{center}

The checklist also says that what is asked in time series and applied deep learning ``completely depends on the projects you put in your resume''.
Use Chapter 30's resume drill for those.
